% Physical p.13 / Roman XI. Source-aligned English translation.
\IllusieSetPhysicalPageSize{450.6bp}{704bp}
\begin{illusiefrontpage}{illusie:I:intro:p013}{}
\begin{center}\fontsize{10}{12}\selectfont XI\end{center}
\vspace{12bp}
\phantomsection\hypertarget{illusie:I:intro:p013:p1}{}%
the usual extension defined by the first-order jets of a module. This invariant
gives rise to a theory of Chern classes for perfect complexes, presented in
Chapter V after some preliminaries on derived categories of filtered complexes.
In particular, the theory includes that of the characteristic polynomial of an
endomorphism of a perfect complex, which resolves a question raised in
(SGA 6 XIV).

\vspace{4bp}
\phantomsection\hypertarget{illusie:I:intro:p013:p2}{}%
Chapter VI\textsuperscript{(*)} prepares the ground for obstruction calculations
concerning deformations of group schemes. It consists, for the most part, of
general calculations on the cohomology of topoi associated with certain diagrams.
These calculations complement those of Deligne--Saint-Donat in (SGA 4 VI),
although our account is logically independent of theirs. The most striking result
is the expression of $\operatorname{Ext}$ groups of modules over a ring as
``spatial'' cohomology groups of a topos defined from Eilenberg--Mac-Lane spectra.
As a by-product, we obtain a remarkable explicit formula for the resolution sought
by Grothendieck in (SGA 7 VII 3.5.4). The idea of the construction, which goes
back to Mac-Lane\textsuperscript{(3)}, was suggested to us by L. Breen, who makes
essential use of this resolution in his recent work on
$\operatorname{Ext}^{1}(G_a,G_a)$\textsuperscript{(4)}.

\vspace{4bp}
\phantomsection\hypertarget{illusie:I:intro:p013:p3}{}%
Chapter VII contains obstruction calculations for first-order deformations of
torsors under a flat group scheme, of noncommutative flat group schemes, and of
flat commutative group schemes with a scheme of operator rings. The obstructions
are again described as cup products with certain fundamental classes, among which
one finds (in the case of torsors) a generalization of the class constructed by
Atiyah in\textsuperscript{(5)}. The results of this

\IllusieNoteRule
\phantomsection\hypertarget{illusie:I:intro:p013:ns}{}%
\IllusieNote{*}{Chapters VI through VIII will be published in a second volume, which will appear shortly.}
\phantomsection\hypertarget{illusie:I:intro:p013:n3}{}%
\IllusieNote{3}{S. Mac-Lane, \emph{Homologie des anneaux et des modules}, Colloque de Topologie algébrique, Louvain, 1956, p. 55--80.}
\phantomsection\hypertarget{illusie:I:intro:p013:n4}{}%
\IllusieNote{4}{L. Breen, work in preparation.}
\phantomsection\hypertarget{illusie:I:intro:p013:n5}{}%
\IllusieNote{5}{M. F. Atiyah, \emph{Analytic connexions on fibre bundles}, Mexico Symposium, 1958.}
\end{illusiefrontpage}
